\answer{ex:22:1}{(a)~$32\%$ e $100\%$. (b)~$R_\uparrow-R_\downarrow\ge2\sqrt{2000\,\text{Hz}/1\,\text{s}}\approx89$~Hz, apenas $4.5\%$ da frequência total.}
\answer{ex:22:2}{(a)~$n\ge4(p_1q_1+p_2q_2)/\Delta^2\approx560$. (b)~$\approx56$.}
\answer{ex:22:3}{(a)~Balanço detalhado: $\rho PK$ é simétrico. (b)~A fração de erros é igual à média harmônica $1/\langle1/P\rangle=0.18$, contra $0.5$ sem realimentação.}
\answer{ex:22:4}{(a)~Um paralelogramo de área $4h^2=0.04$ (levando em conta o limite de $p$, numericamente $0.045$). (b)~$\approx0.18$.}
\answer{ex:22:5}{$0.49$: parte das culturas vai para uma região em que quase sempre erra e vagueia por lá. Com o limite, $1.00$: num conjunto limitado, a densidade $\propto1/P$ se concentra na região absorvente.}
\answer{ex:22:6}{(a)~São precisos $\ge5$ níveis; $8$ eletrodos bastam. (b)~Não: seria preciso $r_{\max}\ge25/T=100$~Hz.}
\answer{ex:22:7}{Problema em aberto. O tempo da busca aleatória cresce com $d$ aproximadamente como a razão de volumes $(L/h)^d$; o da busca com o sinal do erro, linearmente em $d$. Separa as hipóteses um experimento com troca: depois do erro, um estímulo previsível mas forte; depois do acerto, um imprevisível mas fraco; são necessárias algumas dezenas de culturas em cada grupo.}
